\studytypesubsection{LLMs for New Software Engineering Tools}
\label{sec:llms-for-new-software-engineering-tools}

\studytypeparagraph{Description}

LLMs are being integrated into new tools that support software engineers in their daily tasks (e.g., to assist in code comprehension~\cite{DBLP:conf/chi/YanHWH24} and test case generation~\cite{DBLP:journals/tse/SchaferNET24}).
One way of integrating LLM-based tools into software engineers' workflows is using GenAI agents.
Unlike traditional LLM-based tools, these agents are capable of acting autonomously and proactively, are often tailored to meet specific user needs, and can interact with external environments~\cite{DBLP:conf/icse-seip/TakerngsaksiriPTTZJLCCW25,wiesinger2025agents,yang2025code}.
From an architectural perspective, GenAI agents can be implemented in various ways~\cite{wiesinger2025agents}.
However, they generally share three key components: (1) a reasoning mechanism that guides the LLM (often enabled by advanced prompt engineering), (2) a set of tools to interact with external systems (e.g., APIs or databases), and (3) a user communication interface that extends beyond traditional chat-based interactions~\cite{DBLP:conf/icsm/RichardsW24, DBLP:journals/tmlr/SumersYN024, DBLP:journals/corr/abs-2309-07870}.
Researchers can also test and compare different tool architectures to increase artifact quality and developer satisfaction.

\studytypeparagraph{Examples}

\citeauthor{DBLP:conf/chi/YanHWH24} proposed \emph{IVIE}, a tool integrated into the VS Code graphical interface that generates and explains code using LLMs~\cite{DBLP:conf/chi/YanHWH24}.
The authors focused more on the presentation, providing a user-friendly interface to interact with the LLM. 
\citeauthor{DBLP:journals/tse/SchaferNET24}~\cite{DBLP:journals/tse/SchaferNET24} presented a large-scale empirical evaluation on the effectiveness of LLMs for automated unit test generation.
They presented \emph{TestPilot}, a tool that implements an approach in which the LLM is provided with prompts that include the signature and implementation of a function under test, along with usage examples extracted from the documentation.
\citeauthor{DBLP:conf/icsm/RichardsW24} introduced a preliminary GenAI agent designed to assist developers in understanding source code by incorporating a reasoning component grounded in the theory of mind~\cite{DBLP:conf/icsm/RichardsW24}.
